\documentclass[11pt]{letter}
\usepackage[margin=1in]{geometry}
\usepackage{hyperref}

\signature{Francisco Molina Burgos\\Independent Researcher\\yatrogenesis@proton.me}
\address{}

\begin{document}

\begin{letter}{Editorial Office\\Physical Review E\\American Physical Society}

\opening{Dear Editor,}

I am pleased to submit the manuscript entitled \textbf{``Stochastic Informational Primacy in 2D Phase Transitions: A CUSUM Analysis of Persistent Homology''} for consideration in Physical Review E.

\textbf{Summary of the work:} This manuscript investigates the temporal relationship between topological information structure and metric ordering during crystallization. Using molecular dynamics simulations of a 2D Lennard-Jones system, we demonstrate that regime changes in persistence entropy ($S_{H1}$) statistically precede hexatic order consolidation in 73.3\% of independent trials, with a mean lead time of $\sim$750 integration steps.

\textbf{Key novelty and contribution:} Unlike previous studies that treat topological descriptors as passive diagnostics, our work introduces a \textit{predictive} framework where topological instability is detected via industrial-grade change-point algorithms (CUSUM). This methodological advance yields three contributions:

\begin{enumerate}
    \item \textbf{Robust detection:} CUSUM achieves 73.3\% precursor detection vs. 60\% with derivative-based methods, with substantially reduced variance.
    \item \textbf{Mechanistic insight:} Results support ``biased co-emergence'' where informational collapse facilitates (but does not deterministically control) material ordering.
    \item \textbf{Methodological bridge:} We connect TDA techniques with statistical process control, opening pathways for real-time transition monitoring in experimental systems.
\end{enumerate}

\textbf{Scope and audience:} The manuscript addresses fundamental questions in statistical physics regarding information-matter relationships during non-equilibrium transitions. We believe PRE's readership in statistical mechanics, soft matter, and computational physics will find these results relevant.

\textbf{Reproducibility:} All simulation code, analysis pipelines, and raw data are publicly available at \url{https://github.com/Yatrogenesis/TDA-Phase-Transitions}, enabling independent verification.

\textbf{Suggested reviewers:} Given the interdisciplinary nature combining TDA with phase transitions, we suggest reviewers with expertise in (1) persistent homology applications to physical systems, (2) 2D melting/crystallization, or (3) change-point detection in time series.

The manuscript has not been published elsewhere and is not under consideration by another journal. All authors have approved the manuscript for submission.

Thank you for considering this work.

\closing{Sincerely,}

\end{letter}
\end{document}
